\documentclass{report}
\usepackage{amsfonts, amsmath}
\newcommand\R{{\mathbb R}}
\newcommand\Q{{\mathbb Q}}
\newcommand\N{{\mathbb N}}
\newcommand\ve{\varepsilon}
\newcommand\Co{{\mathcal C}}
\pagestyle{empty}
\begin{document}
\large
\centerline{\Huge Fisica Matematica II}
\renewcommand{\thefootnote}{ } 
\footnote{\sl Avete 2 ore di tempo. Ogni esercizio
vale dieci punti. La sufficienza si ottiene con un punteggio $\geq$
18. Solo le risposte chiaramente giustificate saranno prese in
considerazione.}
\renewcommand{\thefootnote}{\arabic{footnote}} 
\vskip.1cm
\centerline{\Large Prima Prova di Autovalutazione}
\vskip1cm
\begin{enumerate}
\item Si risolva il seguente problema di Cauchy
\[
u^2u_x+u_y=0
\]\[
u(x,0)=x.
\]
\item Sia $j:\R_+\to\R_+$, $\operatorname{supp}j\subset [0,1]$ e, per
ogni $x\in\R^n$
\[
j_\ve(x):=\frac{j(\ve^{-1}\|x\|)}{\ve^{n}\int_{\R^n}j(\|y\|)dy} .
\]
Si dimostri che, per ogni $f\in\Co^0(\R^n,\R)$
\[
\lim_{\ve\to 0}\int_{\R^n}j_\ve(x-y)f(y)f(y)dy=f(x).
\]
\item Si risolva il seguente problema di Dirichlet nel dominio
$\Omega:=[-1,1]^2$ 
\[
\Delta u=0
\]\[
u(-1,y)=u(1,y)=0\quad \forall y\in[-1,1]
\]\[
u(x,-1)=u(x,1)=\sin(\pi x)\quad\forall x\in[-1,1].  
\]
\item Si calcoli il Laplaciano in coordinate polari in $\R^3$ e si
trovino le equazioni per la parte radiale e angolare di funzioni armoniche.
\end{enumerate}
\end{document}
